\subsection*{A Full Croqtile-target C++ Code Example}

Here is a complete example that adds two arrays of the same shape and element type. Its three parts are labelled in the listing: a device helper, the tileflow function and the host program that calls it:

\begin{lstlisting}[language=croqtile,caption={A Full Croqtile-target C Code Example}]
// Device Program
__device__ void kernel(int * a, int * b, int * c, int n) {
  for (int i = 0; i < n; ++i) c[i] = a[i] + b[i];
}
// Tileflow Program
__co__ s32 [6, 17, 128] ele_add(s32 [6, 17, 128] lhs, s32 [6, 17, 128] rhs) {
  s32 [lhs.span] output; // Use same shape as lhs
  parallel p by 6 {
    with index in [17, 4] { // Tiling factors
      foreach index {
        lhs_load = dma.copy lhs.chunkat(p, index) => local;
        rhs_load = dma.copy rhs.chunkat(p, index) => local;
        local s32 [lhs_load.span] l1_out;
        // Call kernel with loaded data
        call kernel(lhs_load.data, rhs_load.data, l1_out, |lhs_load.span|);
        // Store result back to output
        dma.copy l1_out => output.chunkat(p, index);
      }
    }
  }
  return output;
}
// Host Program
int main() {
  // Define data arrays
  croqtile::s32 a[6][17][128] = {0};
  croqtile::s32 b[6][17][128] = {0};
  // Fill arrays with data
  std::fill_n(&a[0][0][0], sizeof(a) / sizeof(a[0][0][0]), 1);
  std::fill_n(&b[0][0][0], sizeof(b) / sizeof(b[0][0][0]), 2);
  // Call Croqtile function (data movement and device kernel execution)
  auto res = ele_add(croqtile::make_spanview<3>(&a[0][0][0], {6, 17, 128}),
                     croqtile::make_spanview<3>(&b[0][0][0], {6, 17, 128}));
  // Verification: check correctness of results
  for (size_t i = 0; i < res.shape()[0]; ++i)
    for (size_t j = 0; j < res.shape()[1]; ++j)
      for (size_t k = 0; k < res.shape()[2]; ++k)
        if (a[i][j][k] + b[i][j][k] != res[i][j][k]) {
          std::cerr << "result does not match.\n";
          abort();
        }
  std::cout << "Test Passed\n" << std::endl;
}
\end{lstlisting}

The subsequent sections will explain the different parts of the code.
